% answers_short.tex —— 2005 年·普通高等学校招生考试·广东卷 · 简短答案速查
% 由 tools/make_short_answers.py 自 answers.tex 的 \gksolhead/\gkanswer 提取，
% 并与 2005/2005_数学真题_广东_答案要点.md 逐题比对后生成，勿手改（改源 MD 后重跑）。

\begin{gknotice}
\textbf{使用说明}：本节只给最终答案便于快速核对，完整推导见解析版。答案为编写者独立推导，\textbf{非官方标准答案}。
\end{gknotice}

\gkpart{选择题}
\begin{center}\begin{tabular}{ccccccccccc}
\toprule
题号 & 1 & 2 & 3 & 4 & 5 & 6 & 7 & 8 & 9 & 10 \\
\midrule
答案 & B & D & A & D & B & D & C & C & A & B \\
\bottomrule
\end{tabular}\end{center}

\gkpart{填空题}
\begin{itemize}[leftmargin=2.6em]
  \item[\textbf{11.}] $(-\infty, 0)$
  \item[\textbf{12.}] $4$
  \item[\textbf{13.}] $\pm \dfrac{\sqrt{2}}{2}$
  \item[\textbf{14.}] $5$；$\dfrac{n(n-1)}{2} - 1$
\end{itemize}

\gkpart{解答题}
\begin{itemize}[leftmargin=2.6em]
  \item[\textbf{15.}] 化简结果为 $f(x) = 4\cos 2x$；值域为 $[-4, 4]$，最小正周期为 $\pi$。
  \item[\textbf{16.}] (1) 证明见解析；(2) $\arccos\dfrac{3\sqrt{34}}{34}$
  \item[\textbf{17.}] (1) $y = 3x^2 + \dfrac{2}{3}$；(2) 存在，最小值为 $1$
  \item[\textbf{18.}] (1) $\xi$ 的分布列为：当 $k \in \{0, 1, 2, \dots, n-1\}$ 时，$P(\xi = k) = \left(\dfrac{t}{s+t}\right)^k \dfrac{s}{s+t}$；当 $k = n$ 时，$P(\xi = n) = \left(\dfrac{t}{s+t}\right)^n$；(2) $\mathbf{E}\xi = \dfrac{t}{s}\left[1 - \left(\dfrac{t}{s+t}\right)^n\right]$
  \item[\textbf{19.}] (1) 非奇非偶函数；(2) $802$ 个
  \item[\textbf{20.}] (1) $y = kx + \dfrac{k^2+1}{2}$（$-2 \leqslant k \leqslant 0$）；(2) $2(\sqrt{6} - \sqrt{2})$
\end{itemize}
